\documentclass[10pt,a4paper]{amsart}
\usepackage[margin=19mm]{geometry}
\usepackage{amsmath,amssymb,mathtools}
\usepackage[scr=rsfs,cal=euler]{mathalfa}
\usepackage{xfrac}
\usepackage[unicode,hidelinks,bookmarks=false]{hyperref}
\newcommand{\ie}{i.e.\ }
\newcommand{\cf}{cf.\ }
\newcommand{\resp}{respectively\ }
\newcommand{\Subsection}{\subsection}
\providecommand{\nobreakdash}{\nobreak\hbox{-}}
\setlength{\emergencystretch}{2em}
\multlinegap0pt
\usepackage{bm}
\usepackage{bbm}
\usepackage{dsfont}
\usepackage{xspace}
\usepackage{cancel}
\usepackage{mathtools}
\usepackage{amsthm}
\makeatletter
\@ifundefined{theorem}{\newtheorem{theorem}{Theorem}}{}
\makeatother
\def\tc{{\mathsf{TC}}}
\title[Autonomous motion in changing environment, fibrations and re]{Autonomous motion in changing environment, fibrations and reaction mechanisms}
\author{Michael Farber, Stefan Kurz, Mathias Pillin}
\date{Source: arXiv:2511.11042v2 (CC BY 4.0)}
\begin{document}
\maketitle

\noindent\emph{Source and licence.} Autonomous motion in changing environment, fibrations and reaction mechanisms. Version arXiv:2511.11042v2 (CC BY 4.0), distributed under the Creative Commons Attribution 4.0 International licence, as recorded in the arXiv licence metadata for this exact version and shown on the abstract page https://arxiv.org/abs/2511.11042v2. Full rights notice: licenses/arxiv-2511.11042-CC-BY-4.0.txt.\par\smallskip

\noindent\textit{Editorial scope.} Counted unit: the Theorem whose source label is \texttt{comp} in the article (source0.tex lines 492--495, source label \texttt{comp}), delivered together with the article's own complete original proof of it (source0.tex lines 496--515, one \texttt{proof} environment of 20 lines). No mathematics is written, repaired or re-derived: statement and proof are the article's own text line for line. Required context retained uncounted: pi (lines 162--164); thm:8 (lines 334--339); eq:s (lines 387--389); eq:16d (lines 398--400). Every definition, display and label of those lines is retained unchanged. Omitted from this package and not redistributed: 1--161, 165--333, 340--386, 390--397, 401--491, 516--1385. Nothing inside a retained excerpt is omitted or altered: every retained body is delivered exactly as the article's lines are written, so no omission receipt of any kind accompanies this package. Citations: the retained excerpts carry no citation command at all, so no bibliography entry of the article is transcribed and no reference is redistributed. Reference adaptation: none was needed and none was made; every reference inside the delivered unit resolves inside it, so no unresolved reference or citation is produced. Printed slips kept literal: no printed word, symbol, hypothesis or number of the counted statement or of its proof was altered. Layout and class: the article's own class is \texttt{amsart} with packages including geometry, amssymb, epstopdf, enumerate, xy, amsmath, graphicx, xcolor; the wrapper is the lane's pinned \texttt{amsart} editorial wrapper at 10pt on A4, which restates the theorem-like environments the retained text needs and adds the article's own macro definitions that the retained text needs (\texttt{\textbackslash tc}), with extra wrapper packages (bm, bbm, dsfont, xspace, cancel, mathtools). The delivered numbering is the unit's own. Assets: no retained span contains a figure environment, a graphics-inclusion command, a PDF-page inclusion, a table or a TikZ picture; no image file is copied into this package, the package declares no asset dependency and no image file is redistributed with it. Licence: version arXiv:2511.11042v2 of this article is distributed under the Creative Commons Attribution 4.0 International licence, as recorded in the arXiv licence metadata for this exact version and shown on the abstract page https://arxiv.org/abs/2511.11042v2. A full rights notice is delivered at licenses/arxiv-2511.11042-CC-BY-4.0.txt. Characters: every retained excerpt is pure ASCII, so the delivered engine, XeLaTeX with its default Latin Modern text font, reports no missing character.
\begin{eqnarray}\label{pi}
\Pi: E^I_B \to E^2_B, \quad \mbox{where}\quad \Pi(\gamma) =(\gamma(0), \gamma(1)),
\end{eqnarray}

\begin{theorem}\label{thm:8}
For any locally trivial bundle $p:E\to B$ with Hausdorff and paracompact base $B$ one has the equality
\begin{eqnarray}\label{in:11}
\widetilde\tc[p:E\to B]= \tc[p:E\to B].
\end{eqnarray}  
\end{theorem}

\begin{eqnarray}\label{eq:s}
s: E^2_B\to E^I_B
\end{eqnarray}

\begin{eqnarray}\label{eq:16d}
\tilde s: B^I\times_{B^2}E^2\to E^I
\end{eqnarray}

\begin{theorem}\label{comp}
The extended algorithm $\tilde s$, see (\ref{eq:16d}), has complexity equal to the complexity of the initial algorithm $s$, see (\ref{eq:s}). 
Therefore, under the assumption that the motion of external conditions is known in advance, the motion planning problem in the situation with changing external conditions has complexity equal the complexity of the initial problem with stationary external conditions.
\end{theorem}

\begin{proof}
Suppose that the initial parametrised algorithm $s$ was continuous on an open subset 
$U\subset E^2_B$. Practically this means that the output of the parametrised algorithm $s$ varies continuously with variations of the input 
(i.e. the initial and final configurations and the state of the external conditions) while the input remains in $U$. 
From formula (\ref{eq:16d}) we see that the extended algorithm $\tilde s$ will be continuous on the subset 
$$
\tilde U= Q^{-1}(U)\subset B^I\times_B E^2_B,
$$
where $Q: B^I\times_B E^2_B\to E^2_B$ is given by 
$$Q(\gamma, e, e')=(e, P_\gamma(e')).$$
Here $e, e'\in E$ and $\gamma\in B^I$ satisfy $p(e)=p(e')=\gamma(0)$. Recall that $P_\gamma$ denotes the operator of parallel transport along $\gamma$. 

Let $k=\tc[p:E\to B]$ denote the complexity of the initial problem. Then there exists an open cover 
$E^2_B =U_0\cup U_1\cup\dots \cup U_k$ with a continuous section $s_i:U_i\to E^I_B$ of fibration (\ref{pi}). 
Then the sets $Q^{-1}(U_i)$ form an open cover of the set 
$B^I\times_B E^2_B$ with continuous sections $\tilde s_i$ of the bundle $\tilde \Pi$. This means that the sectional category of the bundle 
$\tilde \Pi$ is smaller than or equal to $k=\tc[p:E\to B]$. On the other hand, the bundle $\tilde \Pi$ contains the bundle $\Pi$ as a restriction of a subset of the constant paths $B\subset B^I$ which implies the opposite inequality. 
We see that the sectional category of the bundle $\tilde \Pi$ equals 
the parametrised topological complexity $\tc[p:E\to B]$. This gives an alternative proof of Theorem \ref{thm:8}. 
\end{proof}

\end{document}
